\documentclass[12pt,a4paper]{article}

% =========================
% Paquetes
% =========================
\usepackage[utf8]{inputenc}
\usepackage[T1]{fontenc}
\usepackage{lmodern}
\usepackage[spanish]{babel}
\usepackage{geometry}
\usepackage{enumitem}
\usepackage{xcolor}
\usepackage{graphicx}
\usepackage[hidelinks]{hyperref}
\usepackage[most]{tcolorbox}
\usepackage{microtype}
\usepackage{titlesec}
\usepackage{fancyhdr}
\usepackage{tabularx}
\usepackage{array}
\usepackage{ragged2e}
\usepackage{tikz}

% =========================
% Configuración de página
% =========================
\geometry{
	margin=2.25cm,
	headheight=16pt
}

% =========================
% Colores institucionales
% Azul eléctrico, amarillo y plomo
% =========================
\definecolor{imegreen}{RGB}{0,102,255}       % Azul eléctrico principal
\definecolor{imegreenDark}{RGB}{0,76,153}    % Azul institucional oscuro
\definecolor{imegreenSoft}{RGB}{230,240,255} % Azul suave
\definecolor{imedark}{RGB}{75,85,99}         % Plomo oscuro
\definecolor{imegray}{RGB}{107,114,128}      % Plomo medio
\definecolor{imelight}{RGB}{245,247,250}     % Fondo plomo claro
\definecolor{imegold}{RGB}{255,204,0}        % Amarillo institucional
\definecolor{imered}{RGB}{185,28,28}

% =========================
% Hipervínculos
% =========================
\hypersetup{
	colorlinks=true,
	linkcolor=imegreenDark,
	urlcolor=imegreenDark,
	citecolor=imegreenDark
}

% =========================
% Listas
% =========================
\setlist[itemize]{
	leftmargin=1.2cm,
	itemsep=0.25em,
	label=\textcolor{imegreen}{\small\textbullet}
}

\setlist[enumerate]{
	leftmargin=1.2cm,
	itemsep=0.25em,
	label=\textcolor{imegreen}{\textbf{\arabic*.}}
}

% =========================
% Encabezado y pie de página
% =========================
\pagestyle{fancy}
\fancyhf{}
\fancyhead[L]{\small\textcolor{imegray}{Bases Intersemestres 2026-II}}
\fancyhead[R]{\small\textcolor{imegray}{Futsal y Vóley Mixto}}
\fancyfoot[C]{\small\textcolor{imegray}{\thepage}}
\renewcommand{\headrulewidth}{0.4pt}
\renewcommand{\footrulewidth}{0pt}

% =========================
% Estilo de títulos
% =========================
\titleformat{\section}
{\Large\bfseries\color{imedark}}
{\colorbox{imegreen}{\textcolor{white}{\thesection}}}
{0.75em}
{}
[\vspace{0.25em}{\color{imegreen}\titlerule[0.8pt]}]

\titleformat{\subsection}
{\large\bfseries\color{imegreenDark}}
{\thesubsection}
{0.65em}
{}

\titlespacing*{\section}{0pt}{1.2em}{0.8em}
\titlespacing*{\subsection}{0pt}{0.9em}{0.45em}

% =========================

% =========================
% Comandos generales
% =========================
\newcommand{\dato}[2]{%
	\noindent\textcolor{imegreenDark}{\textbf{#1:}} #2\par\vspace{0.25em}
}

\newcommand{\datoTabla}[2]{%
	\textbf{\textcolor{imegreenDark}{#1}} & #2\\[0.35em]
}

\newcommand{\logoCarrera}{%
	\IfFileExists{assets/logo-carrera.png}{%
		\includegraphics[width=2.7cm]{assets/logo-carrera.png}%
	}{%
		\fbox{\parbox[c][2.2cm][c]{3.8cm}{\centering Logo de la carrera}}%
	}%
}

% =========================
% Cajas visuales
% =========================
\newtcolorbox{cajaResumen}[1]{
	enhanced,
	breakable,
	colback=white,
	colframe=imegreen,
	coltitle=white,
	colbacktitle=imegreen,
	fonttitle=\bfseries,
	title=#1,
	boxrule=0.9pt,
	arc=3mm,
	left=3mm,
	right=3mm,
	top=2mm,
	bottom=2mm,
	before skip=0.45cm,
	after skip=0.45cm
}

\newtcolorbox{cajaTarjeta}[1]{
	enhanced,
	colback=white,
	colframe=imegreen,
	coltitle=white,
	colbacktitle=imegreen,
	fonttitle=\bfseries,
	title=#1,
	boxrule=0.9pt,
	arc=3mm,
	left=3mm,
	right=3mm,
	top=2mm,
	bottom=2mm
}

\newtcolorbox{cajaImportante}[1]{
	enhanced,
	breakable,
	colback=yellow!12,
	colframe=imegold,
	coltitle=black,
	colbacktitle=imegold,
	fonttitle=\bfseries,
	title=#1,
	boxrule=0.9pt,
	arc=3mm,
	left=3mm,
	right=3mm,
	top=2mm,
	bottom=2mm,
	before skip=0.45cm,
	after skip=0.45cm
}

\newtcolorbox{cajaAdvertencia}[1]{
	enhanced,
	breakable,
	colback=gray!8,
	colframe=imedark,
	coltitle=white,
	colbacktitle=imedark,
	fonttitle=\bfseries,
	title=#1,
	boxrule=0.9pt,
	arc=3mm,
	left=3mm,
	right=3mm,
	top=2mm,
	bottom=2mm,
	before skip=0.45cm,
	after skip=0.45cm
}

\newtcolorbox{cajaPago}{
	enhanced,
	colback=imegreenSoft,
	colframe=imegreen,
	boxrule=1pt,
	arc=4mm,
	left=4mm,
	right=4mm,
	top=4mm,
	bottom=4mm,
	width=0.78\textwidth,
	drop shadow
}

% =========================


% Protocolo de la Mesa de Control: que hacer antes, durante y despues de cada partido.
% Se basa en bases-intersemestres-2026-II.tex; lo que las bases no aclaran va como criterio (C1, C2...)
% y se aprueba en el Acta de criterios del final.

\fancyhead[L]{\small\textcolor{imegray}{Protocolo de la Mesa de Control}}
\fancyhead[R]{\small\textcolor{imegray}{Intersemestres 2026-II}}

% Etiqueta de criterio propuesto (se aprueba en el acta final)
\newcommand{\crit}[1]{\hspace{1pt}\colorbox{imegold}{\scriptsize\bfseries\textcolor{black}{C#1}}}
\newcommand{\casilla}{\fbox{\rule{0pt}{5pt}\hspace{5pt}}}
\newcommand{\linea}[1]{\rule{#1}{0.4pt}}

\begin{document}

% =====================================================================
% PORTADA BREVE
% =====================================================================
\thispagestyle{empty}
\begin{tikzpicture}[remember picture,overlay]
	\fill[imegreen] (current page.north west) rectangle ([yshift=-0.45cm]current page.north east);
	\fill[imegold] ([yshift=-0.45cm]current page.north west) rectangle ([yshift=-0.62cm]current page.north east);
\end{tikzpicture}
\begin{center}
	\vspace*{0.3cm}
	\logoCarrera\par\vspace{0.3cm}
	{\large\textbf{\textcolor{imedark}{Escuela Profesional de Ingeniería Mecánica Eléctrica}}\par}
	\vspace{0.4cm}
	{\Huge\textbf{\textcolor{imedark}{Protocolo de la Mesa de Control}}\par}
	\vspace{0.15cm}
	{\Large\textcolor{imegreenDark}{Gran Campeonato Intersemestres 2026-II · Futsal y Vóley Mixto}\par}
	\vspace{0.15cm}
	{\large Jueves 29 de octubre de 2026 · Complejo Deportivo Pereyra\par}
\end{center}

\begin{cajaResumen}{Para qué sirve este documento}
	Es la guía de trabajo de la mesa: qué revisar \textbf{antes}, qué controlar \textbf{durante} y qué cerrar \textbf{después} de cada partido, y qué decidir en cada situación. Se basa en las Bases oficiales. Donde las bases no dicen nada o se contradicen, se propone un criterio marcado así: \crit{1}. Esos criterios se aprueban y firman en el \textbf{Acta de criterios} (última sección) \textbf{antes del primer partido}; desde ese momento son obligatorios y no se discuten en cancha.
\end{cajaResumen}

\begin{cajaTarjeta}{Formato real tras el sorteo}
	\renewcommand{\arraystretch}{1.25}
	\noindent\begin{tabularx}{\linewidth}{>{\RaggedRight\arraybackslash}p{0.2\linewidth} >{\RaggedRight\arraybackslash}X >{\RaggedRight\arraybackslash}X}
		& \textbf{Futsal} & \textbf{Vóley Mixto}\\ \hline
		Equipos & 15 en 4 grupos (A, B, C de 4; D de 3) & 8 en 2 grupos de 4\\
		Partidos & 29: 21 de grupos, 4 cuartos, 2 semis, 3.\textsuperscript{er} lugar y final & 15: 12 de grupos, 2 semis y final\\
		Llave & C1: 1A--2B · C2: 1B--2A · C3: 1C--2D · C4: 1D--2C & S1: 1A--2B · S2: 1B--2A\\
		Duración & 2 tiempos de 10 min, sin descanso & 1 set a 20 (dif. 2, tope 22); final al mejor de 3 (20, 20, 10)\\
		Canchas & Cancha A y Cancha B & Losa de vóley\\
		Tabla & Victoria 3, empate 1, derrota 0 & Victoria 2, derrota 0\\
	\end{tabularx}
	\smallskip
	{\small Las bases describen el formato con cupo completo (18 y 12 equipos); según la Disposición final 4 se ajustó al número de inscritos. Horarios y orden: los del \textbf{fixture impreso} de este documento.}

\begin{cajaAdvertencia}{Todo se lleva en papel}
	La mesa trabaja solo con este documento impreso: fixture y hoja de resultados, planillas por partido, tablas de posiciones, llaves, registros y bases. No se usa celular, internet ni aplicaciones para decidir. Lo que no está escrito en las planillas no existe para la mesa.
\end{cajaAdvertencia}

\begin{cajaTarjeta}{Contenido de este documento}
	\begin{enumerate}
		\item Protocolo de la mesa y Acta de criterios.
		\item Fixture y hoja de resultados (Futsal y Vóley).
		\item Anexos: planilla en blanco de reserva (A), registro de sanciones (B), hoja de reclamo (C) y registro de caja (D).
		\item Grupos, tablas de posiciones y llaves (Futsal y Vóley).
		\item Planillas de control, una por partido, ya llenadas con número, hora, cancha y equipos.
		\item Bases oficiales completas (sustento de toda decisión).
	\end{enumerate}
\end{cajaTarjeta}
\end{cajaTarjeta}

% =====================================================================
\section{Organización de la mesa}
% =====================================================================
\begin{itemize}
	\item \textbf{Una mesa por disciplina} (futsal y vóley), cada una con:
	\begin{itemize}
		\item \textbf{Jefe de mesa:} decide en el momento, firma el acta y es el único que habla con delegados y árbitro sobre sanciones.
		\item \textbf{Planillero:} llena la planilla del partido, lleva el tiempo de llamado y pasa el resultado a la hoja de resultados y a la tabla de posiciones.
		\item \textbf{Control de identidad:} revisa nóminas firmadas, DNI y fichas de matrícula impresas.
		\item \textbf{Tesorería:} cobra arbitraje y reclamos, entrega recibo y lleva el registro de caja.
	\end{itemize}
	\item La \textbf{Mesa Evaluadora} (jefes de ambas mesas + un miembro de la comisión) resuelve reclamos y todo lo no previsto. Sus decisiones son inapelables.
	\item Las dos mesas designan un \textbf{mensajero} que lleva avisos escritos entre ellas para los \textbf{equipos que juegan en las dos disciplinas} (sección~\ref{sec:cruces}).
	\item \textbf{Material:} este documento impreso, nóminas firmadas y fichas de matrícula de cada equipo, lapiceros, reloj, silbato de respaldo, chalecos, cintas de capitán, talonario de recibos y sencillo.
\end{itemize}

% =====================================================================
\section{Antes del partido}
% =====================================================================

\subsection{Línea de tiempo}
\renewcommand{\arraystretch}{1.3}
\noindent\begin{tabularx}{\linewidth}{>{\bfseries}p{2.2cm} >{\RaggedRight\arraybackslash}X}
	Un partido antes & Ambos equipos deben estar presentes (previsión de W.O.). La mesa los llama por altavoz o en persona.\\
	$-10$ min & El delegado o capitán entrega la \textbf{lista del partido} (titulares y suplentes) elegida solo entre los 10 inscritos. Se revisa identidad y se cobra arbitraje.\\
	$-5$ min & Revisión de uniforme, cinta de capitán y, en vóley, composición en cancha. Sorteo de chalecos si coinciden colores.\\
	Hora programada & Empieza el partido. Si falta un equipo corre la \textbf{tolerancia de 5 minutos} (en vóley, desde el primer llamado del árbitro).\\
	$+5$ min & Si el equipo no completa el mínimo (4 jugadores): \textbf{W.O.}\\
\end{tabularx}

\subsection{Lista del partido e identidad}
\begin{enumerate}
	\item Solo pueden jugar los inscritos en la \textbf{nómina} del equipo (máx. 10). La lista del partido no se modifica una vez iniciado el encuentro.
	\item Cada jugador se identifica con \textbf{DNI o carné} y se compara con la nómina y su ficha de matrícula 2026-II (egresados: documento que acredite que fueron estudiantes de la carrera).
	\item Se revisa el \textbf{registro de sanciones} (Anexo B): quien tenga roja en su partido anterior no juega este partido.
	\item Si hay duda sobre la matrícula, el jugador debe mostrar su \textbf{ficha de matrícula impresa} junto con su DNI. Si no puede comprobarlo en ese momento, \textbf{no juega} hasta comprobarlo \crit{1}.
\end{enumerate}

\begin{cajaImportante}{Jugadores que no cumplen: qué hacer antes del partido}
	\renewcommand{\arraystretch}{1.3}
	\noindent\begin{tabularx}{\linewidth}{>{\RaggedRight\arraybackslash}p{0.45\linewidth} >{\RaggedRight\arraybackslash}X}
		\textbf{Caso} & \textbf{Decisión}\\ \hline
		No pertenece a la carrera, sin ficha válida, con suspensión académica, de otro semestre o con menos de 2 cursos del semestre & Se \textbf{retira} al jugador; no juega. El equipo juega si mantiene el mínimo de 4. No hay más sanción.\\
		No figura en la nómina del equipo & No juega. Si llega a entrar a la cancha es \textbf{suplantación}: el equipo queda \textbf{eliminado} (bases).\\
		Inscrito también en otro equipo de la \textbf{misma disciplina} & No juega por ninguno hasta que la Mesa Evaluadora resuelva. Si llega a jugar por dos equipos, el equipo infractor queda \textbf{eliminado} (bases).\\
		Jugador de futsal que también juega vóley (otra disciplina) & \textbf{Permitido}: la restricción es por disciplina \crit{2}.\\
		Jugadora invitada en vóley (no E.P.I.M.E.) & Permitida: máximo 2, solo mujeres y deben figurar en la nómina.\\
		Suspendido por roja & No juega este partido.\\
	\end{tabularx}
\end{cajaImportante}

\subsection{Uniforme, capitán y colores}
\begin{itemize}
	\item Camisetas de un mismo color. Si un equipo no tiene uniforme, \textbf{el partido se juega igual} y se anota en la planilla; la sanción se aplica al final (sección~\ref{sec:despues}).
	\item El capitán lleva cinta de color distinto a la camiseta y es el único que habla con el árbitro.
	\item Si coinciden los colores, se sortea en presencia de los capitanes: el que pierde usa chalecos o voltea la camiseta.
\end{itemize}

\subsection{Composición en vóley mixto}
\begin{itemize}
	\item En cancha: \textbf{máximo 5 varones} y \textbf{mínimo 1 dama} (6 jugadores).
	\item Equipo sin damas (ni invitadas): juega con \textbf{5 varones}. Si \textbf{ambos} equipos son solo de varones, juegan con 6.
	\item Mínimo para iniciar: 4 jugadores, siempre con un máximo de 5 varones en cancha.
\end{itemize}

\subsection{Arbitraje}
\begin{itemize}
	\item Cada equipo paga en la mesa \textbf{antes de iniciar}, en todas las fases: \textbf{futsal S/ 7.50} y \textbf{vóley S/ 6.00} por equipo (S/ 15 y S/ 12 por partido). Se entrega recibo.
	\item El equipo que no paga \textbf{juega igual}, pero queda como deuda y no puede jugar su siguiente partido hasta pagarla \crit{3}.
\end{itemize}

\subsection{Walkover (W.O.)}
\begin{itemize}
	\item Pasados los 5 minutos de tolerancia, el equipo que no tiene al menos 4 jugadores inscritos y aptos pierde por W.O.
	\item Marcador: futsal \textbf{3--0}. Vóley: \textbf{20--0} (un set); en la final \textbf{2--0 en sets} (20--0, 20--0) \crit{4}.
	\item \textbf{Ambos ausentes:} en grupos, derrota para los dos (0 puntos, 0--3 en contra cada uno en futsal; 0--20 en vóley). En eliminación, avanza el que tenga mejor posición en su tabla de grupo; si es igual, sorteo \crit{5}.
\end{itemize}

\subsection{Checklist de inicio (se marca en la planilla)}
\casilla\ Lista del partido entregada por ambos equipos \quad
\casilla\ Identidad verificada \quad
\casilla\ Sin suspendidos en cancha\\[3pt]
\casilla\ Arbitraje cobrado \quad
\casilla\ Uniformes y capitanes revisados \quad
\casilla\ Composición vóley correcta \quad
\casilla\ Capitanes firman la planilla

% =====================================================================
\section{Durante el partido}
% =====================================================================

\subsection{Qué anota la mesa}
\begin{itemize}
	\item \textbf{Futsal:} goles (número del jugador y minuto), \textbf{faltas acumuladas por tiempo} (avisar al árbitro al llegar a 5: desde la 6.\textsuperscript{a} es tiro libre sin barrera desde 10 m; se reinician en el segundo tiempo), tarjetas y expulsiones (el expulsado no se reemplaza durante 2 minutos o hasta que el rival anote), y cambios (ilimitados, solo entre los de la lista).
	\item \textbf{Vóley:} puntos de cada set, 1 tiempo por set por equipo, cambios (libres) y que después de cada cambio sigan siendo máximo 5 varones en cancha.
	\item El \textbf{tiempo lo controla solo el árbitro}; la mesa lleva un cronómetro de respaldo.
\end{itemize}

\subsection{Situaciones durante el partido}
\renewcommand{\arraystretch}{1.3}
\noindent\begin{tabularx}{\linewidth}{>{\RaggedRight\arraybackslash}p{0.38\linewidth} >{\RaggedRight\arraybackslash}X}
	\textbf{Situación} & \textbf{Decisión de la mesa}\\ \hline
	Se detecta un jugador que \textbf{no pertenece a la carrera} & Se le saca de inmediato y \textbf{el partido continúa} \crit{6}. Al terminar se aplica la sanción: el equipo infractor \textbf{pierde los puntos} de ese partido (sección~\ref{sec:despues}).\\
	Entra un jugador que \textbf{no está en la nómina} & Suplantación: el árbitro detiene el partido y el equipo queda \textbf{eliminado} (bases).\\
	Vóley: quedan 6 varones en cancha o un equipo sin damas juega con 6 & El árbitro detiene el juego hasta corregir; el punto en curso no se repite \crit{7}.\\
	Futsal: una expulsión deja a un equipo con menos de 4 & Partido terminado; gana el rival con al menos \textbf{3--0} (o la diferencia mayor si ya la tenía).\\
	Vóley: un equipo queda con menos de 4 & Partido terminado; gana el rival y el set se completa a su favor hasta 20 (en la final, también los sets restantes) \crit{8}.\\
	Agresión física o verbal directa & Jugador separado de inmediato; con agresión directa el equipo es \textbf{retirado del campeonato}. La mesa redacta el informe.\\
	Hinchas ingresan al campo & Se suspende el partido y gana el rival; si se repite, el equipo es separado del campeonato. Se aplica en \textbf{ambas disciplinas} \crit{9}.\\
	Reclamo en pleno juego & Solo el capitán, al árbitro. No se aceptan reclamos de apreciación arbitral. Lo demás se presenta al terminar.\\
	Lluvia, corte de luz u otra causa de fuerza mayor & Se detiene el reloj. Si se reanuda en menos de 15 min, sigue con el mismo marcador y tiempo. Si no, la Mesa Evaluadora reprograma el tiempo restante en el primer hueco libre \crit{10}.\\
	Daño al recinto o arrojar basura & Se anota en la planilla; el equipo paga el daño y recibe la sanción al final \crit{11}.\\
\end{tabularx}

% =====================================================================
\section{Después del partido}\label{sec:despues}
% =====================================================================

\subsection{Empates}
\begin{itemize}
	\item \textbf{Futsal, fase de grupos:} el empate vale \textbf{1 punto para cada equipo}; no hay penales \crit{12}.
	\item \textbf{Futsal, cuartos, semifinales y 3.\textsuperscript{er} lugar:} \textbf{3 penales por equipo} y luego uno por equipo hasta que haya ganador.
	\item \textbf{Futsal, final:} \textbf{5 penales por equipo} y luego uno por equipo.
	\item Solo patean los jugadores que estaban en cancha al terminar el partido. La mesa anota el orden y el resultado de la tanda.
	\item \textbf{Vóley:} no hay empates (sets con diferencia de 2, tope 22).
\end{itemize}

\subsection{Sanciones que se aplican al cerrar el partido}
\renewcommand{\arraystretch}{1.3}
\noindent\begin{tabularx}{\linewidth}{>{\RaggedRight\arraybackslash}p{0.33\linewidth} >{\RaggedRight\arraybackslash}X >{\RaggedRight\arraybackslash}X}
	\textbf{Caso} & \textbf{Fase de grupos} & \textbf{Eliminación directa}\\ \hline
	Jugó un jugador que \textbf{no pertenece a la carrera} (detectado durante o al término) & El infractor \textbf{pierde los puntos}: 0 puntos y el partido se registra como derrota 0--3 (futsal) / 0--20 (vóley), o el marcador real si el rival ganó por más \crit{13}. & El infractor queda \textbf{eliminado} y avanza el rival \crit{13}.\\
	Equipo \textbf{sin uniforme} & Se le resta \textbf{1 punto en la tabla} \crit{14}. & Se resta 1 gol (futsal) o 1 punto del set (vóley) al marcador final; si así queda empatado, se define por penales (futsal) \crit{14}.\\
	Daño al recinto o basura & Se le resta 1 punto en la tabla y paga el daño \crit{11}. & Paga el daño; amonestación en el acta \crit{11}.\\
\end{tabularx}

\subsection{Acta, reclamos y registro}
\begin{enumerate}
	\item El árbitro y ambos capitanes \textbf{firman la planilla} con el resultado (y penales, si hubo). Firmar es dar conformidad.
	\item \textbf{Reclamos:} el delegado o capitán los escribe al dorso de la planilla (o en el Anexo C) \textbf{dentro de los 5 minutos} después del partido, con el sustento. Se pagan \textbf{S/ 5.00 al presentarlo}; si el reclamo procede, se devuelven \crit{15}. Pasados los 5 minutos no se recibe.
	\item La Mesa Evaluadora resuelve \textbf{antes del siguiente partido de cualquiera de los dos equipos} y comunica la decisión a ambos delegados. Si no alcanza, el siguiente partido se juega y la decisión se aplica después.
	\item El planillero pasa el resultado a la \textbf{hoja de resultados} (marcador, penales o W.O., y la sanción si la hubo) y actualiza la \textbf{tabla de posiciones} del grupo. Las planillas firmadas se guardan en orden de partido: son el sustento ante cualquier reclamo.
	\item Se actualiza el \textbf{registro de sanciones} (Anexo B): roja = no juega el siguiente partido; segunda roja en el campeonato = fuera del campeonato. Dos amarillas en el mismo partido = roja; las amarillas no se arrastran a otros partidos.
\end{enumerate}

\subsection{Cierre de grupos y paso a la llave}
\begin{itemize}
	\item \textbf{Desempate futsal:} 1) puntos, 2) diferencia de goles, 3) goles a favor, 4) menos goles en contra, 5) resultado entre los empatados, 6) sorteo en la mesa con los delegados \crit{16}.
	\item \textbf{Desempate vóley:} 1) puntos, 2) diferencia de puntos, 3) puntos a favor, 4) menos puntos en contra, 5) sorteo (bases).
	\item Al terminar cada grupo, el jefe de mesa firma la tabla, escribe el 1.° y 2.° en la llave y avisa a los delegados.
	\item Los cruces de la llave son los del cuadro de la portada y no se cambian.
\end{itemize}

% =====================================================================
\section{Equipos que juegan futsal y vóley}\label{sec:cruces}
% =====================================================================
Las bases permiten mover un partido solo si coincide con el de vóley del mismo equipo. Procedimiento \crit{17}:
\begin{enumerate}
	\item Antes de empezar el día, ambas mesas marcan en su fixture los partidos de los equipos que juegan las dos disciplinas (con sus horas).
	\item Si coinciden, \textbf{termina primero el partido que ya empezó}; el otro espera al equipo hasta \textbf{10 minutos} después de que termine, sin W.O.
	\item El mensajero lleva el aviso escrito a la otra mesa y ambas anotan la espera en la planilla.
\end{enumerate}

% =====================================================================
\section{Acta de criterios}
% =====================================================================
Criterios que completan lo que las bases no aclaran. Se aprueban o modifican \textbf{antes del primer partido} y desde entonces son obligatorios.

\renewcommand{\arraystretch}{1.35}
{\small
\noindent\begin{tabularx}{\linewidth}{c >{\RaggedRight\arraybackslash}X c c}
	& \textbf{Criterio propuesto} & \textbf{Aprob.} & \textbf{Modif.}\\ \hline
	C1 & Si no se puede comprobar la matrícula en el momento, el jugador no juega hasta comprobarla. & \casilla & \casilla\\
	C2 & Un jugador puede estar en un equipo de futsal y en uno de vóley; la prohibición de dos equipos es dentro de la misma disciplina. & \casilla & \casilla\\
	C3 & Quien no paga arbitraje juega, pero no puede jugar su siguiente partido hasta pagar. & \casilla & \casilla\\
	C4 & W.O. en vóley: 20--0; en la final 2--0 en sets (20--0, 20--0). & \casilla & \casilla\\
	C5 & Ambos ausentes: en grupos derrota para los dos; en eliminación avanza el de mejor posición de grupo o por sorteo. & \casilla & \casilla\\
	C6 & Jugador que no pertenece a la carrera detectado durante el partido: se le saca y el partido sigue; la sanción se aplica al final. & \casilla & \casilla\\
	C7 & Vóley con exceso de varones: se detiene hasta corregir; el punto en curso no se repite. & \casilla & \casilla\\
	C8 & Vóley con menos de 4: gana el rival, completando a 20 el set (y los sets restantes en la final). & \casilla & \casilla\\
	C9 & Ingreso de hinchas al campo: misma sanción en futsal que en vóley. & \casilla & \casilla\\
	C10 & Fuerza mayor: hasta 15 min de espera; luego la Mesa Evaluadora reprograma el tiempo restante. & \casilla & \casilla\\
	C11 & Daño o basura: paga el daño y $-1$ punto en tabla (grupos) o amonestación (eliminación). & \casilla & \casilla\\
	C12 & Empate en grupos de futsal: 1 punto cada uno, sin penales (las bases dicen penales también en grupos; se propone no usarlos para que la tabla sea simple). & \casilla & \casilla\\
	C13 & Jugador que no pertenece a la carrera (durante o al término): en grupos, 0 puntos al infractor y derrota 0--3 / 0--20; en eliminación, el infractor queda fuera. & \casilla & \casilla\\
	C14 & Sin uniforme: $-1$ punto en tabla (grupos); $-1$ gol o punto del marcador en eliminación. & \casilla & \casilla\\
	C15 & Reclamo: S/ 5.00 al presentarlo, se devuelven si procede. & \casilla & \casilla\\
	C16 & Desempate futsal: puntos, DG, GF, menos GC, resultado entre ellos, sorteo. & \casilla & \casilla\\
	C17 & Cruce futsal/vóley: termina el que empezó; el otro espera hasta 10 min sin W.O. & \casilla & \casilla\\
\end{tabularx}
}

\vspace{0.4cm}
Modificaciones: \linea{0.8\linewidth}\\[8pt]
\linea{\linewidth}\\[8pt]
\linea{\linewidth}

\vspace{1.4cm}
\noindent
\noindent\begin{tabularx}{\linewidth}{XXX}
	\centering\linea{4.5cm}\\ Presidente de la Comisión &
	\centering\linea{4.5cm}\\ Jefe de mesa Futsal &
	\centering\linea{4.5cm}\\ Jefe de mesa Vóley\\
\end{tabularx}

% =====================================================================
\clearpage
% Generado con: node --env-file=.env.local scripts/mesa-fixture-tex.mjs
\section*{Fixture y hoja de resultados: Futsal}
\addcontentsline{toc}{section}{Fixture y hoja de resultados: Futsal}
{\small El planillero anota aquí el resultado de cada partido al recibir la planilla firmada. Si hubo penales, anotarlos en Obs. (ej.: pen. 3--2). En la fase final, escribir el nombre del equipo clasificado sobre el cruce en gris.}\par
\subsection*{Fase de grupos}
{\footnotesize\renewcommand{\arraystretch}{1.5}
\noindent\begin{tabularx}{\linewidth}{|c|c|c|c|>{\RaggedLeft\arraybackslash}X|c|>{\RaggedRight\arraybackslash}X|p{1.5cm}|}
\hline
\textbf{N.°} & \textbf{Hora} & \textbf{Can.} & \textbf{Fase} & \textbf{Local} & \textbf{Resultado} & \textbf{Visitante} & \textbf{Obs.}\\ \hline
1 & 08:00 & A & Gr. A & Cabo verde & \casilla\,--\,\casilla & Primer semestre & \\ \hline
2 & 08:00 & B & Gr. B & REAL BARZA-VII A & \casilla\,--\,\casilla & Los Chankas & \\ \hline
3 & 08:25 & A & Gr. C & FC PAPAYITAS & \casilla\,--\,\casilla & ETS & \\ \hline
4 & 08:25 & B & Gr. D & Los Traralelitos FC & \casilla\,--\,\casilla & MEKACHEROS & \\ \hline
5 & 08:50 & A & Gr. A & Nottingham Miedo & \casilla\,--\,\casilla & F.C. Real Polquitos & \\ \hline
6 & 08:50 & B & Gr. B & Antidenuncias FC & \casilla\,--\,\casilla & Terrible Resaca & \\ \hline
7 & 09:15 & A & Gr. C & Sport gusto & \casilla\,--\,\casilla & FC ACHUMI & \\ \hline
8 & 10:00 & A & Gr. D & Salsitas al poder de fuego & \casilla\,--\,\casilla & MEKACHEROS & \\ \hline
9 & 10:00 & B & Gr. A & Cabo verde & \casilla\,--\,\casilla & F.C. Real Polquitos & \\ \hline
10 & 10:25 & A & Gr. B & REAL BARZA-VII A & \casilla\,--\,\casilla & Terrible Resaca & \\ \hline
11 & 10:25 & B & Gr. C & FC PAPAYITAS & \casilla\,--\,\casilla & FC ACHUMI & \\ \hline
12 & 10:50 & A & Gr. A & Primer semestre & \casilla\,--\,\casilla & Nottingham Miedo & \\ \hline
13 & 10:50 & B & Gr. B & Los Chankas & \casilla\,--\,\casilla & Antidenuncias FC & \\ \hline
14 & 11:15 & A & Gr. C & ETS & \casilla\,--\,\casilla & Sport gusto & \\ \hline
15 & 12:00 & A & Gr. D & Salsitas al poder de fuego & \casilla\,--\,\casilla & Los Traralelitos FC & \\ \hline
16 & 12:00 & B & Gr. A & Cabo verde & \casilla\,--\,\casilla & Nottingham Miedo & \\ \hline
17 & 12:25 & A & Gr. B & REAL BARZA-VII A & \casilla\,--\,\casilla & Antidenuncias FC & \\ \hline
18 & 12:25 & B & Gr. C & FC PAPAYITAS & \casilla\,--\,\casilla & Sport gusto & \\ \hline
19 & 12:50 & A & Gr. A & F.C. Real Polquitos & \casilla\,--\,\casilla & Primer semestre & \\ \hline
20 & 12:50 & B & Gr. B & Terrible Resaca & \casilla\,--\,\casilla & Los Chankas & \\ \hline
21 & 13:15 & A & Gr. C & FC ACHUMI & \casilla\,--\,\casilla & ETS & \\ \hline
\end{tabularx}}
\subsection*{Fase final}
{\footnotesize\renewcommand{\arraystretch}{2.6}
\noindent\begin{tabularx}{\linewidth}{|c|c|c|c|>{\RaggedLeft\arraybackslash}X|c|>{\RaggedRight\arraybackslash}X|p{1.5cm}|}
\hline
\textbf{N.°} & \textbf{Hora} & \textbf{Can.} & \textbf{Fase} & \textbf{Local} & \textbf{Resultado} & \textbf{Visitante} & \textbf{Obs.}\\ \hline
22 & 14:00 & A & C1 & \vspace*{\fill}{\scriptsize\color{imegray}1.° Grupo A} & \casilla\,--\,\casilla & \vspace*{\fill}{\scriptsize\color{imegray}2.° Grupo B} & \\ \hline
23 & 14:00 & B & C2 & \vspace*{\fill}{\scriptsize\color{imegray}1.° Grupo B} & \casilla\,--\,\casilla & \vspace*{\fill}{\scriptsize\color{imegray}2.° Grupo A} & \\ \hline
24 & 14:25 & A & C3 & \vspace*{\fill}{\scriptsize\color{imegray}1.° Grupo C} & \casilla\,--\,\casilla & \vspace*{\fill}{\scriptsize\color{imegray}2.° Grupo D} & \\ \hline
25 & 14:25 & B & C4 & \vspace*{\fill}{\scriptsize\color{imegray}1.° Grupo D} & \casilla\,--\,\casilla & \vspace*{\fill}{\scriptsize\color{imegray}2.° Grupo C} & \\ \hline
26 & 15:10 & A & S1 & \vspace*{\fill}{\scriptsize\color{imegray}Ganador C1} & \casilla\,--\,\casilla & \vspace*{\fill}{\scriptsize\color{imegray}Ganador C2} & \\ \hline
27 & 15:10 & B & S2 & \vspace*{\fill}{\scriptsize\color{imegray}Ganador C3} & \casilla\,--\,\casilla & \vspace*{\fill}{\scriptsize\color{imegray}Ganador C4} & \\ \hline
28 & 15:55 & A & 3.er lugar & \vspace*{\fill}{\scriptsize\color{imegray}Perdedor S1} & \casilla\,--\,\casilla & \vspace*{\fill}{\scriptsize\color{imegray}Perdedor S2} & \\ \hline
29 & 16:40 & A & Final & \vspace*{\fill}{\scriptsize\color{imegray}Ganador S1} & \casilla\,--\,\casilla & \vspace*{\fill}{\scriptsize\color{imegray}Ganador S2} & \\ \hline
\end{tabularx}}
\clearpage
\section*{Fixture y hoja de resultados: Vóley Mixto}
\addcontentsline{toc}{section}{Fixture y hoja de resultados: Vóley Mixto}
{\small El planillero anota aquí el resultado de cada partido al recibir la planilla firmada. Resultado en puntos del set (en la final, sets y puntos de cada set). En la fase final, escribir el nombre del equipo clasificado sobre el cruce en gris.}\par
\subsection*{Fase de grupos}
{\footnotesize\renewcommand{\arraystretch}{1.5}
\noindent\begin{tabularx}{\linewidth}{|c|c|c|c|>{\RaggedLeft\arraybackslash}X|c|>{\RaggedRight\arraybackslash}X|p{1.5cm}|}
\hline
\textbf{N.°} & \textbf{Hora} & \textbf{Can.} & \textbf{Fase} & \textbf{Local} & \textbf{Resultado} & \textbf{Visitante} & \textbf{Obs.}\\ \hline
1 & 08:00 & Losa & Gr. A & $\pi$-chu-long & \casilla\,--\,\casilla & Quokkas & \\ \hline
2 & 08:35 & Losa & Gr. B & Pupilos de paredes & \casilla\,--\,\casilla & Antidenuncias FC & \\ \hline
3 & 09:10 & Losa & Gr. A & Los invencibles & \casilla\,--\,\casilla & LAS DIVINAS & \\ \hline
4 & 09:45 & Losa & Gr. B & REYES DE LA POTENCIA & \casilla\,--\,\casilla & Salsitas al poder de fuego & \\ \hline
5 & 10:50 & Losa & Gr. A & $\pi$-chu-long & \casilla\,--\,\casilla & LAS DIVINAS & \\ \hline
6 & 11:25 & Losa & Gr. B & Pupilos de paredes & \casilla\,--\,\casilla & Salsitas al poder de fuego & \\ \hline
7 & 12:00 & Losa & Gr. A & Quokkas & \casilla\,--\,\casilla & Los invencibles & \\ \hline
8 & 12:35 & Losa & Gr. B & Antidenuncias FC & \casilla\,--\,\casilla & REYES DE LA POTENCIA & \\ \hline
9 & 13:40 & Losa & Gr. A & $\pi$-chu-long & \casilla\,--\,\casilla & Los invencibles & \\ \hline
10 & 14:15 & Losa & Gr. B & Pupilos de paredes & \casilla\,--\,\casilla & REYES DE LA POTENCIA & \\ \hline
11 & 14:50 & Losa & Gr. A & LAS DIVINAS & \casilla\,--\,\casilla & Quokkas & \\ \hline
12 & 15:25 & Losa & Gr. B & Salsitas al poder de fuego & \casilla\,--\,\casilla & Antidenuncias FC & \\ \hline
\end{tabularx}}
\subsection*{Fase final}
{\footnotesize\renewcommand{\arraystretch}{2.6}
\noindent\begin{tabularx}{\linewidth}{|c|c|c|c|>{\RaggedLeft\arraybackslash}X|c|>{\RaggedRight\arraybackslash}X|p{1.5cm}|}
\hline
\textbf{N.°} & \textbf{Hora} & \textbf{Can.} & \textbf{Fase} & \textbf{Local} & \textbf{Resultado} & \textbf{Visitante} & \textbf{Obs.}\\ \hline
13 & 16:30 & Losa & S1 & \vspace*{\fill}{\scriptsize\color{imegray}1.° Grupo A} & \casilla\,--\,\casilla & \vspace*{\fill}{\scriptsize\color{imegray}2.° Grupo B} & \\ \hline
14 & 17:05 & Losa & S2 & \vspace*{\fill}{\scriptsize\color{imegray}1.° Grupo B} & \casilla\,--\,\casilla & \vspace*{\fill}{\scriptsize\color{imegray}2.° Grupo A} & \\ \hline
15 & 18:10 & Losa & Final & \vspace*{\fill}{\scriptsize\color{imegray}Ganador S1} & \casilla\,--\,\casilla & \vspace*{\fill}{\scriptsize\color{imegray}Ganador S2} & \\ \hline
\end{tabularx}}
\clearpage


\section*{Anexo A: Planilla de control (en blanco, de reserva)}
\addcontentsline{toc}{section}{Anexo A: Planilla de control (en blanco, de reserva)}
\renewcommand{\arraystretch}{1.6}
\noindent Disciplina: \casilla\ Futsal \casilla\ Vóley \quad Fase: \linea{2.6cm} \quad Partido N.°: \linea{1.2cm}\\
Cancha: \linea{2.4cm} \quad Hora programada: \linea{1.6cm} \quad Hora de inicio: \linea{1.6cm}\\
Árbitro: \linea{6cm}

\vspace{0.3cm}
\noindent\begin{tabularx}{\linewidth}{|>{\RaggedRight\arraybackslash}p{4.2cm}|X|X|}
	\hline
	& \textbf{Equipo local} & \textbf{Equipo visitante}\\ \hline
	Nombre & & \\ \hline
	Lista entregada / identidad OK & \casilla\ / \casilla & \casilla\ / \casilla\\ \hline
	Arbitraje pagado (S/) & & \\ \hline
	Uniforme / cinta capitán & \casilla\ / \casilla & \casilla\ / \casilla\\ \hline
	Vóley: varones en cancha (máx. 5) & & \\ \hline
	Goles o puntos (por set) & & \\ \hline
	Faltas 1.\textsuperscript{er} T / 2.\textsuperscript{o} T (futsal) & & \\ \hline
	Amarillas (N.° jugador) & & \\ \hline
	Rojas (N.° jugador) & & \\ \hline
	Penales & & \\ \hline
	Jugadores retirados (motivo) & & \\ \hline
	Sanciones aplicadas (C\ldots) & & \\ \hline
	\textbf{Resultado final} & & \\ \hline
	Firma capitán & & \\ \hline
\end{tabularx}

\vspace{0.3cm}
\noindent Incidencias: \linea{0.83\linewidth}\\[6pt]
\linea{\linewidth}\\[6pt]
\linea{\linewidth}

\vspace{0.9cm}
\noindent
\noindent\begin{tabularx}{\linewidth}{XX}
	\centering\linea{5cm}\\ Árbitro &
	\centering\linea{5cm}\\ Jefe de mesa\\
\end{tabularx}

% =====================================================================
\clearpage
\section*{Anexo B: Registro de sanciones}
\addcontentsline{toc}{section}{Anexo B: Registro de sanciones}
Se revisa antes de cada partido. Roja = no juega el siguiente partido de su equipo. Segunda roja = fuera del campeonato.

\vspace{0.2cm}
\renewcommand{\arraystretch}{1.9}
\noindent\begin{tabularx}{\linewidth}{|p{2.2cm}|X|X|p{1.6cm}|p{2.6cm}|p{2.2cm}|}
	\hline
	\textbf{Disciplina} & \textbf{Jugador} & \textbf{Equipo} & \textbf{Partido} & \textbf{Sanción} & \textbf{No juega en}\\ \hline
	& & & & & \\ \hline & & & & & \\ \hline & & & & & \\ \hline & & & & & \\ \hline
	& & & & & \\ \hline & & & & & \\ \hline & & & & & \\ \hline & & & & & \\ \hline
	& & & & & \\ \hline & & & & & \\ \hline & & & & & \\ \hline & & & & & \\ \hline
	& & & & & \\ \hline & & & & & \\ \hline & & & & & \\ \hline
\end{tabularx}

% =====================================================================
\clearpage
\section*{Anexo C: Hoja de reclamo}
\addcontentsline{toc}{section}{Anexo C: Hoja de reclamo}
\renewcommand{\arraystretch}{1.6}
\noindent Partido N.°: \linea{1.5cm} \quad Disciplina: \linea{2.5cm} \quad Hora de presentación: \linea{2cm} (máx. 5 min después del partido)\\
Equipo que reclama: \linea{6cm} \quad Contra: \linea{5cm}\\
Presentado por (delegado o capitán): \linea{8.5cm}\\
Pago S/ 5.00: \casilla\ Recibo N.° \linea{2cm}

\vspace{0.3cm}
\noindent\textbf{Motivo del reclamo y sustento:}\\[6pt]
\linea{\linewidth}\\[10pt]\linea{\linewidth}\\[10pt]\linea{\linewidth}\\[10pt]\linea{\linewidth}\\[10pt]\linea{\linewidth}

\vspace{0.5cm}
\noindent\textbf{Resolución de la Mesa Evaluadora:} \casilla\ Procede (se devuelven S/ 5.00) \quad \casilla\ No procede\\[6pt]
\linea{\linewidth}\\[10pt]\linea{\linewidth}\\[10pt]\linea{\linewidth}

\vspace{1.2cm}
\noindent
\noindent\begin{tabularx}{\linewidth}{XXX}
	\centering\linea{4.5cm}\\ Reclamante &
	\centering\linea{4.5cm}\\ Delegado rival &
	\centering\linea{4.5cm}\\ Mesa Evaluadora\\
\end{tabularx}

% =====================================================================
\clearpage
\section*{Anexo D: Registro de caja}
\addcontentsline{toc}{section}{Anexo D: Registro de caja}
Cada cobro (arbitraje o reclamo) se anota con su recibo. Al cerrar el día, el total debe coincidir con el efectivo.

\vspace{0.2cm}
\renewcommand{\arraystretch}{1.75}
\noindent\begin{tabularx}{\linewidth}{|p{1.4cm}|p{1.2cm}|X|p{2.4cm}|p{1.5cm}|p{1.6cm}|}
	\hline
	\textbf{Recibo} & \textbf{Partido} & \textbf{Equipo} & \textbf{Concepto} & \textbf{Monto} & \textbf{Firma}\\ \hline
	& & & & & \\ \hline & & & & & \\ \hline & & & & & \\ \hline & & & & & \\ \hline
	& & & & & \\ \hline & & & & & \\ \hline & & & & & \\ \hline & & & & & \\ \hline
	& & & & & \\ \hline & & & & & \\ \hline & & & & & \\ \hline & & & & & \\ \hline
	& & & & & \\ \hline & & & & & \\ \hline & & & & & \\ \hline & & & & & \\ \hline
	& & & & & \\ \hline & & & & & \\ \hline & & & & & \\ \hline & & & & & \\ \hline
	\multicolumn{4}{|r|}{\textbf{Total del día}} & & \\ \hline
\end{tabularx}

\vspace{1.2cm}
\noindent\begin{minipage}{0.45\linewidth}\centering\linea{5cm}\\ Tesorería\end{minipage}\hfill
\begin{minipage}{0.45\linewidth}\centering\linea{5cm}\\ Jefe de mesa\end{minipage}

\end{document}
